\section{Bab VI: Hak dan Kewajiban Warga Negara dalam UUD NRI Tahun 1945}
\subsection{A. Konsep Dasar \& Prinsip}
\begin{itemize}
    \item \textbf{Hak Konstitusional:} Hak yang dijamin dalam UUD NRI Tahun 1945.
    \item \textbf{Kewajiban:} Sesuatu yang harus dilaksanakan dengan penuh tanggung jawab.
    \item \textbf{Prinsip Hubungan (Kausalitas):} Bersifat timbal balik (\textit{reciprocity}) dan tidak dapat dipisahkan; pelaksanaannya harus seimbang.
\end{itemize}

\subsection{B. Rincian Jaminan Hak dan Kewajiban Berdasarkan Pasal}
\begin{description}
    \item[Hukum \& Pemerintahan (Pasal 27):] Kedudukan yang sama di dalam hukum/pemerintahan (Ayat 1); Hak atas pekerjaan dan penghidupan yang layak (Ayat 2); Hak dan kewajiban bela negara (Ayat 3).
    \item[Hak Asasi Manusia (Pasal 28 \& 28A--28J):] Penjabaran HAM secara terinci:
    \begin{itemize}
        \item \textbf{Pasal 28A--28C:} Hak hidup, berkeluarga, perlindungan anak, serta mengembangkan diri lewat pendidikan dan IPTEK.
        \item \textbf{Pasal 28D--28F:} Kepastian hukum yang adil, hak bekerja, kesempatan sama dalam pemerintahan, kebebasan beragama, dan hak berkomunikasi.
        \item \textbf{Pasal 28G--28H:} Perlindungan diri, bebas dari penyiksaan, hak lingkungan hidup sehat dan jaminan sosial.
        \item \textbf{Pasal 28I:} Hak-hak yang tidak dapat dikurangi dalam keadaan apa pun (\textit{non-derogable rights}) dan perlindungan dari diskriminasi.
        \item \textbf{Pasal 28J:} Kewajiban menghormati HAM orang lain dan tunduk pada pembatasan UU.
    \end{itemize}
    \item[Agama \& Kepercayaan (Pasal 29):] Negara berdasar atas Ketuhanan Yang Maha Esa dan jaminan kemerdekaan memeluk agama/beribadat.
    \item[Pertahanan \& Keamanan (Pasal 30):] Hak dan kewajiban warga negara ikut serta dalam usaha pertahanan dan keamanan negara.
    \item[Pendidikan \& Kebudayaan (Pasal 31 \& 32):] Hak mendapat pendidikan (wajib pendidikan dasar dengan anggaran minimal 20\%) dan jaminan pemajuan kebudayaan nasional serta bahasa daerah.
    \item[Perekonomian \& Kesejahteraan (Pasal 33 \& 34):] Perekonomian disusun atas asas kekeluargaan, cabang produksi vital dikuasai negara, serta pemeliharaan fakir miskin dan anak terlantar oleh negara.
\end{description}

\subsection{C. Contoh Penerapan Sehari-hari}
\begin{itemize}
    \item \textbf{Contoh Hak:} Memeluk agama, menyuarakan pendapat, mendapat pendidikan, memiliki properti legal, dan diperlakukan sama di mata hukum.
    \item \textbf{Contoh Kewajiban:} Membayar pajak, taat aturan hukum, menghormati hak orang lain, menjaga sarana umum, dan menjaga kebersihan lingkungan.
\end{itemize}